% Options for packages loaded elsewhere
\PassOptionsToPackage{unicode}{hyperref}
\PassOptionsToPackage{hyphens}{url}
\PassOptionsToPackage{dvipsnames,svgnames,x11names}{xcolor}
\PassOptionsToPackage{space}{xeCJK}
\documentclass[
]{article}
\usepackage{xcolor}
\usepackage[margin=2.5cm]{geometry}
\usepackage{amsmath,amssymb}
\setcounter{secnumdepth}{-\maxdimen} % remove section numbering
\usepackage{iftex}
\ifPDFTeX
  \usepackage[T1]{fontenc}
  \usepackage[utf8]{inputenc}
  \usepackage{textcomp} % provide euro and other symbols
\else % if luatex or xetex
  \usepackage{unicode-math} % this also loads fontspec
  \defaultfontfeatures{Scale=MatchLowercase}
  \defaultfontfeatures[\rmfamily]{Ligatures=TeX,Scale=1}
\fi
\usepackage{lmodern}
\ifPDFTeX\else
  % xetex/luatex font selection
  \setmainfont[]{DejaVu Sans}
  \setmonofont[]{DejaVu Sans Mono}
  \ifXeTeX
    \usepackage{xeCJK}
    \setCJKmainfont[]{Noto Sans CJK SC}
  \fi
  \ifLuaTeX
    \usepackage[]{luatexja-fontspec}
    \setmainjfont[]{Noto Sans CJK SC}
  \fi
\fi
% Use upquote if available, for straight quotes in verbatim environments
\IfFileExists{upquote.sty}{\usepackage{upquote}}{}
\IfFileExists{microtype.sty}{% use microtype if available
  \usepackage[]{microtype}
  \UseMicrotypeSet[protrusion]{basicmath} % disable protrusion for tt fonts
}{}
\makeatletter
\@ifundefined{KOMAClassName}{% if non-KOMA class
  \IfFileExists{parskip.sty}{%
    \usepackage{parskip}
  }{% else
    \setlength{\parindent}{0pt}
    \setlength{\parskip}{6pt plus 2pt minus 1pt}}
}{% if KOMA class
  \KOMAoptions{parskip=half}}
\makeatother
\usepackage{longtable,booktabs,array}
\usepackage{caption}
\captionsetup[table]{skip=6pt}
\newcounter{none} % for unnumbered tables
\usepackage{calc} % for calculating minipage widths
% Correct order of tables after \paragraph or \subparagraph
\usepackage{etoolbox}
\makeatletter
\patchcmd\longtable{\par}{\if@noskipsec\mbox{}\fi\par}{}{}
\makeatother
% Allow footnotes in longtable head/foot
\IfFileExists{footnotehyper.sty}{\usepackage{footnotehyper}}{\usepackage{footnote}}
\makesavenoteenv{longtable}
\setlength{\emergencystretch}{3em} % prevent overfull lines
\providecommand{\tightlist}{%
  \setlength{\itemsep}{0pt}\setlength{\parskip}{0pt}}
\usepackage{bookmark}
\IfFileExists{xurl.sty}{\usepackage{xurl}}{} % add URL line breaks if available
\urlstyle{same}
% fallback for those not using the hyperref driver hyperxmp:
\makeatletter
\@ifundefined{xmpquote}{\newcommand{\xmpquote}[1]{#1}}{}
\makeatother
\hypersetup{
  colorlinks=true,
  linkcolor={Maroon},
  filecolor={Maroon},
  citecolor={Blue},
  urlcolor={Blue},
  pdfcreator={LaTeX via pandoc}}

\author{}
\date{}

\begin{document}

\section{筑基篇课后习题 XII：BSD 猜想的 pat
重新观测}\label{ux7b51ux57faux7bc7ux8bfeux540eux4e60ux9898-xiibsd-ux731cux60f3ux7684-pat-ux91cdux65b0ux89c2ux6d4b}

\begin{quote}
\textbf{声明 (Declaration)}: 本文\textbf{不代表任何数学结论}
(non-mathematical-claim), 仅为基于 Lean 4
形式化的一种\textbf{实验观测报告} (experimental observation report)。

{\def\LTcaptype{none} % do not increment counter
\begin{longtable}[]{@{}
  >{\raggedright\arraybackslash}p{(\linewidth - 6\tabcolsep) * \real{0.2500}}
  >{\raggedright\arraybackslash}p{(\linewidth - 6\tabcolsep) * \real{0.2500}}
  >{\raggedright\arraybackslash}p{(\linewidth - 6\tabcolsep) * \real{0.2500}}
  >{\raggedright\arraybackslash}p{(\linewidth - 6\tabcolsep) * \real{0.2500}}@{}}
\toprule\noalign{}
\begin{minipage}[b]{\linewidth}\raggedright
习题
\end{minipage} & \begin{minipage}[b]{\linewidth}\raggedright
主题
\end{minipage} & \begin{minipage}[b]{\linewidth}\raggedright
Lean 形式化
\end{minipage} & \begin{minipage}[b]{\linewidth}\raggedright
DOI
\end{minipage} \\
\midrule\noalign{}
\endhead
\bottomrule\noalign{}
\endlastfoot
exercise\_i\_12 & BSD 猜想的 pat 重新观测（R169） &
PatBSDConjecture.lean & 10.5281/zenodo.21917249 \\
\end{longtable}
}
\end{quote}

\textbf{Exercise XII for the Foundation-Building Chapter: The BSD
Conjecture Revisited from the PAT Perspective}

\begin{quote}
习题定位：筑基篇课后习题系列第十二题。用 mechanics-pat-observation skill
观测 BSD 猜想（Birch and Swinnerton-Dyer，Clay
千禧年问题之一）。全部新定理只锚筑基篇定理（R085/R136②③/R138/R146/R159/R161），不用外部引理。
\end{quote}

\emph{2026-08-13 · Internal research exercise · Lean 4 / mathlib v4.32.2
· 5 new theorems PROVED no sorry · 算器神魂论筑基篇课后习题 XII }

\begin{center}\rule{0.5\linewidth}{0.5pt}\end{center}

\subsection{习题陈述}\label{ux4e60ux9898ux9648ux8ff0}

\begin{quote}
对 BSD 猜想（椭圆曲线 E 的 L 函数在 s=1 的零点阶数 = 莫德尔-威尔群 E(ℚ)
的秩）做 pat 重新观测。唯一论点：\textbf{L(E,s) = 素数相位乘积（Euler
积）；莫德尔-威尔秩 = 独立互锁对数；解析秩（零点阶数）= 折叠类深度；BSD
= 两个计数一致。}
\end{quote}

\begin{center}\rule{0.5\linewidth}{0.5pt}\end{center}

\subsection{零、总论：BSD 的 pat
转译}\label{ux96f6ux603bux8bbabsd-ux7684-pat-ux8f6cux8bd1}

BSD 猜想经典陈述：ord\_\{s=1\} L(E,s) = rank E(ℚ)------解析秩 = 代数秩。

\textbf{pat 结构对应}（mechanics-pat-observation skill）：

{\def\LTcaptype{none} % do not increment counter
\begin{longtable}[]{@{}
  >{\raggedright\arraybackslash}p{(\linewidth - 4\tabcolsep) * \real{0.3333}}
  >{\raggedright\arraybackslash}p{(\linewidth - 4\tabcolsep) * \real{0.3333}}
  >{\raggedright\arraybackslash}p{(\linewidth - 4\tabcolsep) * \real{0.3333}}@{}}
\toprule\noalign{}
\begin{minipage}[b]{\linewidth}\raggedright
BSD 结构
\end{minipage} & \begin{minipage}[b]{\linewidth}\raggedright
pat 结构
\end{minipage} & \begin{minipage}[b]{\linewidth}\raggedright
锚定
\end{minipage} \\
\midrule\noalign{}
\endhead
\bottomrule\noalign{}
\endlastfoot
L(E,s) = ∏\_p L\_p(s)（Euler 积） & 素数相位乘积（素数 = 方向 log p） &
R159/R146 \\
E(ℚ) = ℤ\^{}r ⊕ tors（秩 r） & r 对独立互锁（生成元 = 互锁对） &
R161/R136②③ \\
ord\_\{s=1\} L(E,s)（解析秩） & 折叠类深度（零点 = 折叠类） &
R085/R138 \\
★BSD：解析秩 = 代数秩 & 折叠类深度 = 互锁对数（两个计数一致） &
CONJECTURE \\
\end{longtable}
}

\begin{center}\rule{0.5\linewidth}{0.5pt}\end{center}

\subsection{一、素数 = 方向 log p（L
函数素数因子结构）}\label{ux4e00ux7d20ux6570-ux65b9ux5411-log-pl-ux51fdux6570ux7d20ux6570ux56e0ux5b50ux7ed3ux6784}

\textbf{★核心：L(E,s) 的每个素数因子 = 一个素数方向的相位结构（log p
单相位）。}

\subsubsection{论证}\label{ux8bbaux8bc1}

\begin{enumerate}
\def\labelenumi{\arabic{enumi}.}
\tightlist
\item
  \textbf{L(E,s) = Euler 积}：L(E, s) = ∏\_p L\_p(s)，每个素数 p
  贡献一个因子（含 a\_p, p\^{}\{-s\}, p\^{}\{1-2s\}）------与 ζ
  的欧拉乘积同构（C025 zeta\_euler\_product）。
\item
  \textbf{素数 = 方向 log p}（R159 prime\_log\_monophase /
  R146）：log(p\^{}k) = k·log p------素数幂链沿 log 方向 =
  单相位等差链。
\item
  \textbf{L 函数的素数因子 = 素数方向的相位结构}------L(E,s) =
  素数相位的乘积。
\end{enumerate}

\subsubsection{形式化（PatBSDConjecture.lean，1
定理）}\label{ux5f62ux5f0fux5316patbsdconjecture.lean1-ux5b9aux7406}

\begin{itemize}
\tightlist
\item
  \texttt{prime\_direction\_log}：log(p\^{}k) = k·log p------素数 = 方向
  log p（L 函数素数因子结构）。
\end{itemize}

\textbf{验证}：0 sorry。锚定 R159 prime\_log\_monophase。

\begin{center}\rule{0.5\linewidth}{0.5pt}\end{center}

\subsection{二、秩 r = r
对独立互锁（莫德尔-威尔群生成元）}\label{ux4e8cux79e9-r-r-ux5bf9ux72ecux7acbux4e92ux9501ux83abux5fb7ux5c14-ux5a01ux5c14ux7fa4ux751fux6210ux5143}

\textbf{★核心：E(ℚ) = ℤ\^{}r ⊕ tors 的秩 r = 独立生成元数 = r
对独立互锁。}

\subsubsection{论证}\label{ux8bbaux8bc1-1}

\begin{enumerate}
\def\labelenumi{\arabic{enumi}.}
\tightlist
\item
  \textbf{莫德尔-威尔定理}：E(ℚ) 有限生成阿贝尔群 E(ℚ) = ℤ\^{}r ⊕
  tors，秩 r = 独立生成元数。
\item
  \textbf{r 个独立生成元 = r 对独立互锁}（R161
  k\_pairs\_independent\_interlock）：任意 r 对独立互锁全部自洽（R136
  ②③：方向成对声明）。
\item
  \textbf{代数秩 = 互锁对数}------代数侧的结构。
\end{enumerate}

\subsubsection{形式化（PatBSDConjecture.lean，1
定理）}\label{ux5f62ux5f0fux5316patbsdconjecture.lean1-ux5b9aux7406-1}

\begin{itemize}
\tightlist
\item
  \texttt{rank\_interlock\_pairs}：任意 r 对独立互锁自洽------代数秩 =
  互锁对数。
\end{itemize}

\textbf{验证}：0 sorry。锚定 R161。

\begin{center}\rule{0.5\linewidth}{0.5pt}\end{center}

\subsection{三、L 在 s=1 折叠 =
相位锁定（解析秩的折叠类结构）}\label{ux4e09l-ux5728-s1-ux6298ux53e0-ux76f8ux4f4dux9501ux5b9aux89e3ux6790ux79e9ux7684ux6298ux53e0ux7c7bux7ed3ux6784}

\textbf{★核心：L(E,s) 在 s=1 的零点 = 折叠类（R085），解析秩 =
折叠类深度。}

\subsubsection{论证}\label{ux8bbaux8bc1-2}

\begin{enumerate}
\def\labelenumi{\arabic{enumi}.}
\tightlist
\item
  \textbf{零点 = 折叠类}（R085 zero\_is\_fold\_class）：0 =
  折叠类中心------L 在 s=1 的零点落在折叠类。
\item
  \textbf{相位锁定}（R138）：零点 = 相位关系锁定点。
\item
  \textbf{解析秩 = 折叠类深度}：L 在 s=1 的零点阶数 =
  折叠类处的锁定深度。
\end{enumerate}

\subsubsection{形式化（PatBSDConjecture.lean，1
定理）}\label{ux5f62ux5f0fux5316patbsdconjecture.lean1-ux5b9aux7406-2}

\begin{itemize}
\tightlist
\item
  \texttt{lfunction\_fold\_at\_one}：-(t) = 0 ⟹ t = 0------0 =
  折叠类（解析侧的折叠类结构）。
\end{itemize}

\textbf{验证}：0 sorry。锚定 R085 zero\_is\_fold\_class 同型 +
linarith。

\begin{center}\rule{0.5\linewidth}{0.5pt}\end{center}

\subsection{四、★BSD =
两个计数一致（CONJECTURE）}\label{ux56dbbsd-ux4e24ux4e2aux8ba1ux6570ux4e00ux81f4conjecture}

\textbf{★核心：BSD 猜想 = 折叠类深度（解析秩）=
互锁对数（代数秩）------两个计数一致。}

\subsubsection{论证}\label{ux8bbaux8bc1-3}

\begin{enumerate}
\def\labelenumi{\arabic{enumi}.}
\tightlist
\item
  \textbf{代数侧}：秩 r = r 对独立互锁（R161）。
\item
  \textbf{解析侧}：零点 = 折叠类（R085/R138）。
\item
  \textbf{★BSD 转译}：解析秩（L 在 s=1 零点阶数）=
  代数秩（莫德尔-威尔群生成元数）------折叠类深度 = 互锁对数。
\item
  \textbf{诚实边界}：千禧年问题，未解------CONJECTURE（结构转译，非证明）。
\end{enumerate}

\subsubsection{形式化（PatBSDConjecture.lean，2
定理）}\label{ux5f62ux5f0fux5316patbsdconjecture.lean2-ux5b9aux7406}

\begin{itemize}
\tightlist
\item
  \texttt{bsd\_rank\_equality}：★r 对独立互锁 ∧ 0 =
  折叠类------两个计数一致（CONJECTURE）。
\item
  \texttt{bsd\_pat\_perspective}：★全景------素数方向 ∧ 秩 = 互锁对数 ∧
  零点 = 折叠类。
\end{itemize}

\textbf{验证}：0 sorry。合取项分别锚 R159/R161/R085。

\begin{center}\rule{0.5\linewidth}{0.5pt}\end{center}

\subsection{五、习题解答总结}\label{ux4e94ux4e60ux9898ux89e3ux7b54ux603bux7ed3}

{\def\LTcaptype{none} % do not increment counter
\begin{longtable}[]{@{}
  >{\raggedright\arraybackslash}p{(\linewidth - 6\tabcolsep) * \real{0.2500}}
  >{\raggedright\arraybackslash}p{(\linewidth - 6\tabcolsep) * \real{0.2500}}
  >{\raggedright\arraybackslash}p{(\linewidth - 6\tabcolsep) * \real{0.2500}}
  >{\raggedright\arraybackslash}p{(\linewidth - 6\tabcolsep) * \real{0.2500}}@{}}
\toprule\noalign{}
\begin{minipage}[b]{\linewidth}\raggedright
论点
\end{minipage} & \begin{minipage}[b]{\linewidth}\raggedright
内容
\end{minipage} & \begin{minipage}[b]{\linewidth}\raggedright
锚定
\end{minipage} & \begin{minipage}[b]{\linewidth}\raggedright
标签
\end{minipage} \\
\midrule\noalign{}
\endhead
\bottomrule\noalign{}
\endlastfoot
L(E,s) = 素数相位乘积 & 素数 = 方向 log p（Euler 积因子） & R159/R146 &
PROVED \\
秩 = 互锁对数 & r 个生成元 = r 对独立互锁 & R161/R136②③ & PROVED \\
零点 = 折叠类 & L 在 s=1 零点 = 折叠类（R085） & R085/R138 & PROVED \\
★BSD = 计数一致 & 折叠类深度 = 互锁对数 & 上述全部 & CONJECTURE \\
\end{longtable}
}

\textbf{习题的回答（一句话）}：\textbf{BSD 猜想在 pat 视角 =
两个计数一致------解析秩（L 在 s=1 的零点阶数，折叠类深度
R085）等于代数秩（莫德尔-威尔群 E(ℚ) 的生成元数，互锁对数 R161）；L(E,s)
是素数相位乘积（Euler 积），秩是独立互锁对。千禧年未解，CONJECTURE。}

\textbf{诚实边界}：BSD
未解（千禧年问题）------本习题交付结构转译（素数相位乘积 / 秩 = 互锁对数
/ 零点 = 折叠类），CONJECTURE 标注，非证明。

\begin{center}\rule{0.5\linewidth}{0.5pt}\end{center}

\subsection{六、相关对照}\label{ux516dux76f8ux5173ux5bf9ux7167}

{\def\LTcaptype{none} % do not increment counter
\begin{longtable}[]{@{}
  >{\raggedright\arraybackslash}p{(\linewidth - 4\tabcolsep) * \real{0.3333}}
  >{\raggedright\arraybackslash}p{(\linewidth - 4\tabcolsep) * \real{0.3333}}
  >{\raggedright\arraybackslash}p{(\linewidth - 4\tabcolsep) * \real{0.3333}}@{}}
\toprule\noalign{}
\begin{minipage}[b]{\linewidth}\raggedright
文献
\end{minipage} & \begin{minipage}[b]{\linewidth}\raggedright
对应内容
\end{minipage} & \begin{minipage}[b]{\linewidth}\raggedright
备注
\end{minipage} \\
\midrule\noalign{}
\endhead
\bottomrule\noalign{}
\endlastfoot
\textbf{Birch, B., Swinnerton-Dyer, P.} 1965. \emph{Notes on Elliptic
Curves II} & BSD 猜想原始陈述 & pat 转译对象 \\
\textbf{Mordell, L. J.} 1922 / \textbf{Weil, A.} 1928 &
莫德尔-威尔定理（有限生成） & 代数秩的基础 \\
\textbf{R085/R136②③/R138/R146/R159/R161}（筑基篇） &
折叠类/成对/相位锁定/素数/互锁对 & 本框架定理 \\
\end{longtable}
}

\begin{center}\rule{0.5\linewidth}{0.5pt}\end{center}

\emph{筑基篇课后习题 XII · 2026-08-13 · Lean 4 / mathlib v4.32.2 · 5
theorems PROVED no sorry（PatBSDConjecture.lean）· 配套 Lean
形式化：formal/Formal/Toolkit/PatBSDConjecture.lean（快照见
../lean/PatBSDConjecture.lean）}

\end{document}
